\begin{abstract}
Conformal anomaly detection (CAD) provides finite-sample guarantees under exchangeability, but its effectiveness depends on resolving sufficiently small tail probabilities. When the calibration set is small or downstream error control is conservative, the minimum attainable conformal $p$-value can be too large for powerful discovery, even when anomaly scores separate well. This limitation arises in standard CAD, but becomes sharper under covariate shift, where adaptive importance weighting concentrates mass on relevant calibration samples and test points to maintain validity. With that, the standard pipeline encounters what we term \textit{resolution collapse}, where discrete $p$-values fail to clear decision thresholds, while theoretical randomized variants avoid discreteness only by injecting conditional variance. Through mathematical analysis, simulations, and empirical evaluations, we show that these failures are governed by the effective sample size of the calibration set and the conservativeness of error control. Our results suggest that validity alone may not always be enough: effective sample size determines whether conformal guarantees translate into actionable discoveries at the required error-control level.
\end{abstract}

\section{Introduction}\label{sec:introduction}

Anomaly detection aims to identify observations that either deviate from the majority of observations or otherwise do not \textit{conform} to an expected state of normality, indicating a distinct underlying data-generating process \cite{Hawkins1980}. Standard approaches for detection typically rely on heuristic scoring functions that lack finite-sample and non-parametric guarantees regarding the empirical error rates of their decisions, which is problematic in safety-critical applications. 

Conformal anomaly detection (CAD) addresses this by offering a distribution-free and finite-sample valid framework to calibrate heuristic anomaly scores into $p$-values, enabling downstream statistical error control \cite{Bates2023}. CAD calibrates a detector's outputs by comparing the anomaly score of a test point against those computed on held-out calibration samples. Specifically, a test point's $p$-value is its relative score rank among the calibration scores. For this comparison to be valid, calibration samples and non-anomalous test samples must be exchangeable. In dynamic environments where feature distributions shift over time, this assumption is violated. The weighted conformal approach adapts to this covariate shift by assigning higher importance to calibration samples resembling current test instances via likelihood ratios. This localization of calibration weight induces a critical trade-off. As weights concentrate on a smaller effective sample size, empirical $p$-values can become discretely coarse, leading to \textit{resolution collapse}, where the minimum attainable $p$-value fails to meet the discovery threshold required for statistical error control.

Because conformal inference is inherently frequentist, this resolution collapse is not a structural failure of the framework, but rather a manifestation of its statistical transparency. Absent parametric assumptions or subjective priors, CAD refuses to manufacture certainty from effective data scarcity, preferring informative silence (i.e. zero rejections) over ungrounded discoveries. However, because actual error control in CAD operates at the complex interface of conformal calibration and downstream multiple-testing procedures, these fundamental limitations are easily obscured in practice. The subtle interactions between sample size, weight localization under shift, and the targeted error-control level often make it difficult to diagnose a pipeline's degraded statistical power.

We resolve this obscurity by formalizing the mathematical boundaries of the CAD pipeline under shift. We derive the exact limits governing resolution collapse, demonstrate the hidden costs of randomized workarounds, and introduce a pre-computation diagnostic—the detectability ratio—to evaluate a pipeline's viability before evaluating test data.

%While the randomized weighted conformal construction of \cite{Jin2025} replaces the discrete weighted rank by a randomized one, we show that this theoretical fix can come at a practical cost in weighted regimes. When a test point is assigned a large weight (common under shift), smoothing introduces non-negligible auxiliary randomness into the p-value. This randomization improves granularity, but it can also mask weak signals and reduce statistical power while having secondary effects on error control. Consequently, practitioners may be left with an unfavourable trade-off: discrete weighted p-values may have low power because their minimum attainable p-values are inflated by large test weights, whereas randomized weighted p-values can lose efficiency because smoothing noise becomes substantial.

%\paragraph{Contribution.}
%We establish finite-sample detectability limits for CAD: conformal validity does not ensure that \(p\)-values have sufficient tail resolution for powerful testing. We formalize this limit through two failure modes: resolution collapse in discrete conformal \(p\)-values and variance inflation in randomized variants. We diagnose how effective calibration size, multiple testing, and weight localization under shift interact to eliminate power in common scenarios.

\section{Prior Work}\label{sec:prior_works}

CAD \cite{Bates2023,Laxhammar2010} is derived from the field of conformal prediction (CP) \cite{Vovk2005,Shafer2008}, which calibrates classifiers and regressors to produce prediction sets or intervals with coverage guarantees under data exchangeability. CP under covariate shift was established by \citet{Tibshirani2019}, who introduced the weighted exchangeability framework to correct for covariate shift using importance weights, given by \textit{Radon--Nikod{\'y}m} derivatives. For CAD under covariate shift, \citet{Jin2025} extended this framework to multiple testing via \textit{Weighted Conformalized Selection} (WCS). This extension is motivated by the fact that weighted conformal $p$-values need not satisfy Positive Regression Dependence on a Subset (PRDS), a dependence condition under which procedures such as Benjamini--Hochberg (BH) guarantee control of the expected proportion of false discoveries among all discoveries, i.e., False Discovery Rate (FDR) control \cite{Hochberg1995}.

\section{Preliminaries}\label{sec:preliminaries}

We consider unsupervised anomaly detection via one-class classification. Let $\mathcal{D}_{\mathrm{cal}}={z_1,\dots,z_N}$ be a calibration set of inlier observations drawn from a reference distribution $P$. We evaluate a test batch $Z_{\mathrm{test}}={z_{N+1},\dots,z_{N+m}}$, whose points may consist of both inliers and anomalies. For each test point $z_j$, $j\in{N+1,\ldots,N+m}$, we test the null hypothesis $H_{0,j}$ that $z_j$ is an inlier.

Under $H_{0,j}$, the test point $z_j$ is drawn from an inlier test distribution $Q$, which may differ from $P$. The unweighted conformal setting corresponds to $Q=P$, in which case the calibration scores and the score of an inlier test point are exchangeable. Under covariate shift, we assume absolute continuity of the inlier test distribution with respect to the calibration distribution, $Q\ll P$, so that the likelihood ratio $dQ/dP$ is well-defined. Equivalently, regions with positive probability under $Q$ must also be represented under $P$ up to null sets. Let $s:\mathcal{Z}\to\mathbb{R}$ be a score function fixed independently of $\mathcal{D}_{\mathrm{cal}}$ and $Z_{\mathrm{test}}$, with larger values indicating greater deviation from normality. We write $s_i=s(z_i)$.

\subsection{Weighted conformal $p$-values}

Weighted conformal inference accounts for the difference between the calibration distribution $P$ and the inlier test distribution $Q$ through the likelihood ratio
\begin{equation*}
w(z)=\frac{\mathrm{d}Q}{\mathrm{d}P}(z).
\end{equation*}
Since $w$ is generally unknown, practical implementations use an estimate $\widehat w$, typically obtained from a probabilistic domain classifier trained to distinguish samples from the reference distribution $P$ and the target inlier distribution $Q$.

For a test point $z_j$, let $s_j=s(z_j)$, $\omega_j=\omega(z_j)$, and \mbox{$\omega_i=\omega(z_i)$} for $i\in{1,\dots,N}$. Define
\begin{equation*}
A_j^{>}
=
\sum_{i=1}^N \omega_i\mathbb{I}(s_i>s_j),
\qquad
A_j^{=}
=
\sum_{i=1}^N \omega_i\mathbb{I}(s_i=s_j),
\end{equation*}
and let
\begin{equation*}
D_j=\sum_{i=1}^N\omega_i+\omega_j.
\end{equation*}
The standard conformal deterministic and discrete weighted conformal $p$-value \cite{Jin2025} is computed as
\begin{equation*}\label{eq:weighted_ecdf_cons}
p_j^{\mathrm{disc}}
=
\frac{A_j^{>}+A_j^{=}+\omega_j}{D_j}.
\end{equation*}
This is the weighted upper-tail rank of $s_j$, with the test point's own mass included in the normalization. In the oracle case $\omega=w$, under $H_{0,j}$ and the covariate-shift assumption,
\begin{equation*}
\mathbb{P}\left(p_j^{\mathrm{disc}}\le \alpha\right)\le \alpha,
\qquad \forall \alpha\in[0,1].
\end{equation*}
With estimated weights, exact finite-sample validity is not automatic as the procedure targets the oracle construction, and its accuracy depends on the quality of $\widehat w$.

The full test-point mass $\omega_j$ makes $p_j^{\mathrm{disc}}$ conservative and prevents it from reaching zero as
\begin{equation}
p_j^{\mathrm{disc}}\ge \frac{\omega_j}{D_j},
\end{equation}
with equality when $s_j$ exceeds all calibration scores. Setting $\omega\equiv 1$ recovers the usual unweighted conformal upper-tail $p$-value \cite{Bates2023}.

For completeness, we also record the randomized weighted conformal $p$-value \cite{Jin2025}. This variant replaces the deterministic inclusion of the weighted atom at the test score by a uniform randomization over that atom. Let $(U_j\sim\mathrm{Unif}[0,1])$, independently of the data. Then \begin{equation}\label{eq:weighted_ecdf_rand} p_j^{\mathrm{rand}} =
\frac{A_j^{>}+U_j(\omega_j+A_j^{=})}{D_j}.
\end{equation}
Here, the randomization acts on the mass assigned to the test point together with any calibration observations tied with it. If the score distribution is continuous, then $(A_j^{=}=0)$ almost surely, so the only randomized component is the test-point mass. We use this construction primarily as a theoretical extension of the deterministic weighted $p$-value, aligning the notation with weighted conformalized selection, where such randomized $p$-values facilitate finite-sample arguments under weighted exchangeability. Unless stated otherwise, the empirical procedures below use the discrete and deterministic weighted conformal $p$-value.

\subsection{False Discovery Rate Control}

Let
\(\mathcal{J}=\{N+1,\ldots,N+m\}\) denote the test indices and
\(\mathcal{H}_0=\{j\in\mathcal{J}:H_{0,j}\text{ is true}\}\) the true nulls.
For a rejection set \(\mathcal{R}\subseteq\mathcal{J}\),
\[
    \operatorname{FDP}(\mathcal{R})
    =
    \frac{|\mathcal{R}\cap\mathcal{H}_0|}{1\vee|\mathcal{R}|},
    \qquad
    \operatorname{FDR}
    =
    \mathbb{E}\!\left[\operatorname{FDP}(\mathcal{R})\right].
\]
The goal is to construct \(\mathcal{R}\) such that
\(\operatorname{FDR}\le\alpha\).

\paragraph{Benjamini--Hochberg Procedure.}
Given \(p\)-values \(p_{N+1},\ldots,p_{N+m}\), the
BH procedure \cite{Hochberg1995} orders them as
\(p_{(1)}\le\cdots\le p_{(m)}\) and sets
\[
    \widehat{k}
    =
    \max\left\{
        k\in\{1,\ldots,m\}:
        p_{(k)}\le \frac{\alpha k}{m}
    \right\},
\]
with \(\widehat{k}=0\) if the set is empty. It rejects
\[
    \mathcal{R}_{\mathrm{BH}}
    =
    \left\{
        j\in\mathcal{J}:
        p_j\le \frac{\alpha\widehat{k}}{m}
    \right\}.
\]
In the unweighted conformal case, the resulting \(p\)-values are marginally
super-uniform and satisfy a favorable positive-dependence structure, so BH
yields finite-sample FDR control under the standard exchangeability
assumptions \cite{Bates2023}. In the weighted case, marginal validity alone
is insufficient as importance weighting changes the joint dependence among the
test \(p\)-values, and direct BH is no longer justified by the standard BH
argument.

\paragraph{Weighted Conformalized Selection.}
Weighted Conformalized Selection (WCS) \cite{Jin2025} provides a
finite-sample FDR-control layer for weighted conformal \(p\)-values. In what
follows, write \(p_j=p_j^{\mathrm{disc}}\). For each \(j\in\mathcal{J}\), WCS treats \(z_j\) as an additional
weighted calibration point. For \(\ell\in\mathcal{J}\setminus\{j\}\), define
\[
    A_{\ell}^{>,(j)}
    =
    \sum_{i=1}^{N}\omega_i\mathbb{I}\{s_i>s_\ell\}
    +
    \omega_j\mathbb{I}\{s_j>s_\ell\}.
\]
Then
\[
    p_\ell^{(j)}
    =
    \frac{A_{\ell}^{>,(j)}}{D_j}.
\]
WCS then applies BH at level \(\alpha\) to
\[
    \{p_\ell^{(j)}:\ell\in\mathcal{J}\setminus\{j\}\}\cup\{0\},
\]
that is, with the \(j\)-th \(p\)-value forced to zero. Let \(\widehat{\mathcal{R}}_{j\to0}\) denote the rejection set obtained by
applying BH at level \(\alpha\) to
\[
    \{p_\ell^{(j)}:\ell\in\mathcal{J}\setminus\{j\}\}\cup\{0\}.
\]
Set
\[
    a_j=|\widehat{\mathcal{R}}_{j\to0}|,
    \qquad
    t_j=\frac{\alpha a_j}{m}.
\]
Let
\[
    a_j=|\widehat{\mathcal{R}}_{j\to0}|,
    \qquad
    t_j=\frac{\alpha a_j}{m}.
\]
The first-step WCS candidates are
\[
    \mathcal{R}^{(1)}
    =
    \{j\in\mathcal{J}:p_j\le t_j\}.
\]

The candidate set \(\mathcal{R}^{(1)}\) is not used directly. Each threshold
\(t_j\) depends on a hypothesis-specific leave-one-out BH problem, so the
counts \(a_j\) need not agree with the final number of discoveries. This
creates a mismatch between the candidate-wise rejection thresholds and the
denominator in the FDP. WCS therefore applies a pruning step to enforce
self-consistency.

Define
\[
    e_j
    =
    \mathbb{I}\!\left\{
        p_j\le \frac{\alpha a_j}{m}
    \right\},
    \qquad
    B_\eta(r)
    =
    \sum_{j\in\mathcal{J}}
    \mathbb{I}\{e_j=1,\ \eta_j a_j\le r\}.
\]
The final rejection size is
\[
    r_\eta^*
    =
    \max\{r\in\{0,\ldots,m\}:B_\eta(r)\ge r\},
\]
and the WCS rejection set is
\[
    \mathcal{R}_\eta
    =
    \{j\in\mathcal{J}:e_j=1,\ \eta_j a_j\le r_\eta^*\}.
\]
Thus pruning keeps only those first-step candidates whose leave-one-out
rejection count \(a_j\), possibly randomized through \(\eta_j\), is compatible
with the final rejection size. For heterogeneous pruning, let
\(\xi_1,\ldots,\xi_m\stackrel{\mathrm{i.i.d.}}{\sim}\operatorname{Unif}[0,1]\);
for homogeneous pruning, let \(\xi\sim\operatorname{Unif}[0,1]\). The three
variants are
\[
    \eta_j=
    \begin{cases}
        \xi_j, & \text{heterogeneous},\\
        \xi,   & \text{homogeneous},\\
        1,     & \text{deterministic}.
    \end{cases}
\]
Heterogeneous pruning randomizes the self-consistency condition separately
for each candidate, homogeneous pruning uses one shared randomization
variable, and deterministic pruning is fully reproducible but usually more
conservative. In particular, deterministic pruning requires
\(a_j\le r_\eta^*\), whereas the randomized variants may retain candidates
with larger \(a_j\) when the random scaling is favorable.

With oracle weights \(w=\mathrm{d}Q/\mathrm{d}P\), WCS controls the FDR in
finite samples under the weighted covariate-shift assumptions. With estimated
weights, the guarantee depends on density-ratio accuracy. A convenient measure
is
\[
    \widehat{\gamma}
    =
    \sup_{z\in\mathcal{Z}}
    \max\left\{
        \frac{\widehat w(z)}{w(z)},
        \frac{w(z)}{\widehat w(z)}
    \right\}.
\]
Under the independence conditions required for the plug-in analysis, WCS
satisfies the inflated bound
\[
    \operatorname{FDR}
    \le
    \alpha\,
    \mathbb{E}\!\left[
        \frac{\widehat{\gamma}^{\,2}}
        {1+\alpha(\widehat{\gamma}^{\,2}-1)/m}
    \right].
\]
Hence \(\widehat{\gamma}=1\) recovers the oracle guarantee, while larger
values quantify possible FDR inflation due to density-ratio estimation error.

\section{Finite-Sample Power Bounds}
\label{sec:constraints}

In general, CAD describes a pipeline that moves from heuristic anomaly scores to calibrated $p$-values, and finally to a downstream error-control decision that relies on a alignment between the resolution of the calibration step and the stringency of the decision threshold, as a result of the particular error control procedure. Rather than a flaw in the underlying framework, certain setups can expose a structural power boundary. This constraint is inherent to the discrete nature of conformal $p$-values. While present in standard unweighted CAD under small sample sizes, it can become highly pronounced under covariate shift due to weight localization.

\subsection{The Resolution Constraint}
For the standard deterministic weighted conformal construction in Eq.~\eqref{eq:weighted_ecdf_cons}, the smallest attainable $p$-value is strictly lower-bounded by the test point's relative weight:
\begin{equation}\label{eq:resolution_floor}
p^{\mathrm{disc}}_{j} \;\ge\; \frac{\omega_j}{\sum_{i=1}^N \omega_i + \omega_j} = \frac{\omega_j}{W_{\mathrm{total}}}.
\end{equation}
This lower bound establishes an intrinsic \textit{resolution floor}. In a standard unweighted setting ($\omega \equiv 1$), this floor simplifies to a uniform $1/(N+1)$. Even without distribution shift, a small calibration set paired with a strict error-control target can cause this floor to exceed the discovery threshold. Under covariate shift, importance weights localize mass onto a fraction of relevant samples, artificially inflating $\omega_j / W_{\mathrm{total}}$ and lifting this floor for test points in sparse regions.

%\begin{definition}[Detectability Ratio]
%Let $\omega_j/W_{\mathrm{total}}$ be the minimum attainable $p$-value for a test point $j$. For a downstream Benjamini--Hochberg step-up procedure at nominal level $\alpha$ with $r$ putative discoveries, the Detectability Ratio $\delta_j(r)$ is defined as...
%\end{definition}

\paragraph{Detectability Ratio.}
To formalize how this circumstance limits discovery power, consider
a Benjamini--Hochberg step-up procedure applied to a batch of \(m\) tests at
nominal FDR level \(\alpha\). If the final BH rejection count is \(r\), then
the largest rejected \(p\)-value must be no larger than
\((r/m)\alpha\). Suppose the smallest attainable \(p\)-value for test \(j\)
under the discrete calibration scheme is
\(\omega_j/W_{\mathrm{total}}\). We define the \textit{Detectability Ratio}
as
\begin{equation}\label{eq:detectability_ratio}
\delta_j(r)
=
\frac{\displaystyle \omega_j/W_{\mathrm{total}}}
{\displaystyle (r/m)\alpha}.
\end{equation}
If \(\delta_j(r)>1\), then even the most favorable possible score for test
\(j\) cannot yield a \(p\)-value small enough for \(j\) to be included among
\(r\) BH discoveries. In this regime, the calibration grid is too coarse
relative to the downstream BH cutoff. We refer to this failure mode as
\textit{resolution collapse}.

\subsubsection{The Variance Trade-off}
The randomized conformal variant in Equation~\eqref{eq:weighted_ecdf_rand} is often leveraged in theoretical contexts to bypass this resolution floor by smoothing the discrete step at the test point, spreading its probability mass uniformly over an interval. Let $W_{=}(s_j)=\sum_{i=1}^N \omega_i\,\mathbb{I}(s_i=s_j)$. Conditional on the realized calibration data and test scores, $p_j^{\mathrm{rand}}$ is drawn uniformly from an interval of the deterministic width 
\begin{equation}
\mathrm{width}_j = \frac{\omega_j + W_{=}(s_j)}{W_{\mathrm{total}}}.
\end{equation}
In the extreme right-tail case ($s_j > \max_i s_i$), where no ties occur, this simplifies to $p_j^{\mathrm{rand}} \sim \mathrm{Unif}[0, \omega_j/W_{\mathrm{total}}]$.

While this mathematical artifice technically eliminates the resolution floor, it introduces conditional variability that scales quadratically with the test point's weight fraction as
\begin{equation}
\mathrm{Var}\!\left(p_j^{\mathrm{rand}} \mid \mathcal{D}_{\mathrm{cal}}, s_j\right) = \frac{1}{12}\left(\frac{\omega_j}{W_{\mathrm{total}}}\right)^2.
\end{equation}

%\begin{proposition}[Conditional Variance Inflation]
%Conditional on the realized calibration data $\mathcal{D}_{\mathrm{cal}}$ and test score $s_j$, the auxiliary variance introduced by the randomized conformal $p$-value $p_j^{\mathrm{rand}}$ scales quadratically with the test point's relative weight fraction:
%\end{proposition}

In setups where the effective calibration size is restricted, this \textit{variance inflation} acts as auxiliary noise injected into the decision phase. As a result, genuine anomalies may be masked by an unfavorable uniform draw, shifting the failure mode from a deterministic block to a statistical power penalty.

\subsection{The Operational Envelope}
The boundary where these constraints begin to limit pipeline power is defined by a multi-dimensional parameter space. Rather than a systemic failure, power loss is a conditional outcome governed by the interplay of four distinct dimensions:

\begin{enumerate}
    \item \textbf{Effective Calibration Size ($N_{\mathrm{eff}}$):} Following \citet{Kish1965}, weight concentration is measured by $N_{\mathrm{eff}} = (\sum \omega_i)^2 / \sum \omega_i^2$. In standard CAD, $N_{\mathrm{eff}} = N$. Under shift, a lower $N_{\mathrm{eff}}$ increases the test point's relative mass $\omega_j/W_{\mathrm{total}}$, lifting the resolution floor or expanding the randomization interval width.
    \item \textbf{Test Batch Size ($m$):} The size of the evaluated test batch scales the multiple-testing penalty. As $m$ increases relative to the number of true anomalies, the required discovery threshold $(r/m)\alpha$ compresses, sharpening the resolution requirements for calibration.
    \item \textbf{Base Anomaly Rate:} The prevalence of true anomalies dictates the viable range for the rejection count $r$. In low-anomaly environments where $r$ is forced to be small, the denominator of Equation~\eqref{eq:detectability_ratio} shrinks, making the pipeline more sensitive to its resolution limits.
    \item \textbf{Nominal Significance Level ($\alpha$):} Choosing a highly conservative error-control target (low $\alpha$) compresses the decision threshold, requiring lower tail probabilities from the calibration phase.
\end{enumerate}